\chapter{Fuel Cycle and Tritium Breeding}

\section{A Fuel That Does Not Exist}

Every fusion reactor pursuing the deuterium-tritium reaction, the reaction chosen for its comparatively low ignition temperature and the reason nearly every serious fusion program has converged on it, faces a problem that has no analogue in conventional power generation. Deuterium is abundant, extractable from ordinary seawater in effectively unlimited quantity. Tritium is not. It is radioactive, with a half-life of about twelve years, and it does not occur in nature in any quantity worth extracting. Every gram of tritium a fusion reactor consumes must be produced, and the only practical way to produce it at the scale a power plant requires is to produce it inside the reactor itself, as the plant runs.

This single fact reorganizes much of what a fusion reactor has to be. It is not simply a device that burns a fuel supplied from outside, the way a coal plant burns coal delivered by rail. It is a device that must manufacture a meaningful fraction of its own fuel supply, continuously, as a condition of continuing to operate at all. A reactor that achieves excellent plasma performance while consuming tritium faster than it can produce it is not viable in any sense that matters, regardless of its physics, because it is depleting a resource that cannot be replenished from any external source at the necessary scale. Viability, in the fuel cycle, is not a metaphor. It is a literal question of whether the plant can keep feeding itself.

\section{The Breeding Blanket}

Tritium is produced by surrounding the plasma with a blanket containing lithium, positioned to absorb the neutrons that the deuterium-tritium reaction releases in abundance. When a neutron strikes a lithium nucleus, it can produce tritium directly, along with helium and additional energy that contributes to the plant's overall heat output. This is, in a real sense, an elegant arrangement: the same reaction that consumes tritium also produces the neutrons needed to breed more of it, closing a loop between fuel consumption and fuel production within the same physical process.

The elegance, however, is conditional, and the condition is demanding. The ratio of tritium atoms produced to tritium atoms consumed, the breeding ratio, must exceed one by a sufficient margin, not merely on average across the plant's lifetime, but reliably, accounting for tritium lost to radioactive decay during storage, tritium that escapes containment through diffusion into structural materials, tritium consumed in start-up and shutdown transients, and tritium that must be held in reserve to buffer against interruptions in the breeding process itself. A breeding ratio that looks adequate on paper, calculated from an idealized neutron transport model, can prove inadequate in practice once every one of these losses is accounted for, and unlike almost anything else in the reactor, there is no way to import more tritium from outside if the margin turns out to be too thin. The world's existing tritium supply, produced as a byproduct of certain fission reactors, is small, is shrinking as those reactors are decommissioned, and was never intended to fuel commercial fusion power plants at scale.

This is, in the language developed earlier in this book, an admissibility condition of an unusually unforgiving kind. Most of the reactor's admissible region has some margin for temporary excursions, moments where a subsystem drifts close to its boundary and is pulled back by repair processes before real damage occurs. The tritium inventory has much less of this margin, because falling meaningfully behind on breeding is not a transient excursion that can be corrected on a short timescale. It is a slow, compounding deficit that, once established, may take years to recover from, if it can be recovered from at all without an external tritium supply that, at commercial scale, does not currently exist.

\section{Extraction and Recirculation}

Breeding tritium inside the blanket is only the first half of the problem. The tritium then has to be extracted from the blanket material, purified, and recirculated back into the plasma as fuel, fast enough that the fuel cycle does not become a bottleneck on the reactor's operation independent of anything happening in the plasma itself. This extraction and recirculation system is itself a substantial piece of engineering, involving high-temperature chemistry, permeation barriers to prevent tritium from diffusing uncontrolled through structural materials into places it should not go, and continuous monitoring, since tritium is both a valuable fuel and a radiological hazard that must be accounted for precisely at every stage of the cycle.

The viability question this system raises is not simply whether extraction works, but whether it works fast enough and reliably enough to keep pace with the plasma's fuel demand without requiring the plant to hold an impractically large tritium inventory as a buffer against extraction delays. A large buffer inventory sounds like a safety margin, and in one sense it is, but tritium inventory carries its own cost and its own risk: it represents fuel that has been bred but is not yet doing useful work, it decays continuously whether in use or in storage, and a larger standing inventory means a larger quantity of radioactive material that must be safely contained, accounted for, and protected, which is itself a regulatory and security burden with its own viability implications, discussed further in Part VI.

\section{Coupling to the Rest of the Reactor}

The fuel cycle does not operate independently of the reactor's other subsystems, and treating it as though it did is a version of the same mistake Chapter 2 warned against generally. The breeding blanket sits directly behind the first wall, sharing much of the same thermal and neutron environment that determines the first wall's own degradation, a subject taken up in the next chapter. The blanket's structural materials must survive the same punishing neutron flux that makes breeding possible in the first place, and their degradation over time can change the blanket's breeding efficiency, meaning the breeding ratio is not a fixed design parameter but a quantity that itself drifts as the reactor ages. The cooling system that removes heat from the blanket, extracting useful thermal energy for the plant's overall power output, must be designed jointly with the breeding chemistry, since many candidate coolants interact with tritium extraction in ways that constrain both systems simultaneously.

This means the fuel cycle cannot be optimized in isolation any more than plasma, confinement, and heating could be in the previous chapter. A blanket design that maximizes breeding ratio under fresh, undamaged conditions but degrades rapidly under neutron bombardment may fail to sustain an adequate breeding ratio precisely when the plant has been running long enough to have consumed a meaningful fraction of its initial tritium inventory, which is exactly the wrong time for breeding performance to decline. Viability requires evaluating breeding performance across the blanket's full service life, under the same aging conditions the rest of the reactor experiences, not merely at the moment of initial commissioning.

\section{What Self-Sufficiency Actually Requires}

Put together, fuel self-sufficiency in a deuterium-tritium fusion reactor is not a single design target to be hit once. It is a continuously maintained condition, requiring a breeding ratio with real margin above the bare minimum, an extraction and recirculation system fast and reliable enough to avoid bottlenecking plasma operation, a tritium inventory management strategy that balances operational buffer against radiological and security burden, and a blanket design whose breeding performance holds up under the same degradation the rest of the reactor's structural materials experience over decades of neutron exposure.

No existing fusion reactor design has yet been operated long enough, at commercial scale, to demonstrate that all of these conditions can be satisfied jointly across a multi-decade operating life. This is, arguably, the single largest unresolved viability question in fusion engineering today, larger in practical terms than any remaining question in plasma confinement, because unlike confinement, which has been demonstrated to work at increasingly favorable parameters across many experimental devices, sustained tritium self-sufficiency at commercial scale has not yet been demonstrated anywhere, because no device has yet operated long enough, continuously enough, to test it. The next chapter turns to the structural consequences of the same neutron flux that makes breeding possible, and the two problems, breeding tritium and surviving the neutrons that breed it, turn out to be inseparable at the level of reactor design, however separately they are usually discussed.
